\documentclass[12pt,letterpaper]{article}
\usepackage{graphicx,textcomp}
\usepackage{natbib}
\usepackage{setspace}
\usepackage{fullpage}
\usepackage{color}
\usepackage[reqno]{amsmath}
\usepackage{amsthm}
\usepackage{fancyvrb}
\usepackage{amssymb,enumerate}
\usepackage[all]{xy}
\usepackage{endnotes}
\usepackage{lscape}
\newtheorem{com}{Comment}
\usepackage{float}
\usepackage{hyperref}
\newtheorem{lem} {Lemma}
\newtheorem{prop}{Proposition}
\newtheorem{thm}{Theorem}
\newtheorem{defn}{Definition}
\newtheorem{cor}{Corollary}
\newtheorem{obs}{Observation}
\usepackage[compact]{titlesec}
\usepackage{dcolumn}
\usepackage{tikz}
\usetikzlibrary{arrows}
\usepackage{multirow}
\usepackage{xcolor}
\newcolumntype{.}{D{.}{.}{-1}}
\newcolumntype{d}[1]{D{.}{.}{#1}}
\definecolor{light-gray}{gray}{0.65}
\usepackage{url}
\usepackage{listings}
\usepackage{color}

\definecolor{codegreen}{rgb}{0,0.6,0}
\definecolor{codegray}{rgb}{0.5,0.5,0.5}
\definecolor{codepurple}{rgb}{0.58,0,0.82}
\definecolor{backcolour}{rgb}{0.95,0.95,0.92}

\lstdefinestyle{mystyle}{
	backgroundcolor=\color{backcolour},   
	commentstyle=\color{codegreen},
	keywordstyle=\color{magenta},
	numberstyle=\tiny\color{codegray},
	stringstyle=\color{codepurple},
	basicstyle=\footnotesize,
	breakatwhitespace=false,         
	breaklines=true,                 
	captionpos=b,                    
	keepspaces=true,                 
	numbers=left,                    
	numbersep=5pt,                  
	showspaces=false,                
	showstringspaces=false,
	showtabs=false,                  
	tabsize=2
}
\lstset{style=mystyle}
\newcommand{\Sref}[1]{Section~\ref{#1}}
\newtheorem{hyp}{Hypothesis}


\title{Problem Set 4}
\date{Due: November 27, 2025}
\author{Applied Stats/Quant Methods 1}


\begin{document}
	\maketitle
	\section*{Instructions}
	\begin{itemize}
		\item Please show your work! You may lose points by simply writing in the answer. If the problem requires you to execute commands in \texttt{R}, please include the code you used to get your answers. Please also include the \texttt{.R} file that contains your code. If you are not sure if work needs to be shown for a particular problem, please ask.
		\item Your homework should be submitted electronically on GitHub.
		\item This problem set is due before 23:59 on Thursday November 27, 2025. No late assignments will be accepted.
	\end{itemize}



	\vspace{.5cm}
\section*{Question 1: Economics}
\vspace{.25cm}
\noindent 	
In this question, use the \texttt{prestige} dataset in the \texttt{car} library. First, run the following commands:

\begin{verbatim}
install.packages(car)
library(car)
data(Prestige)
help(Prestige)
\end{verbatim} 


\noindent We would like to study whether individuals with higher levels of income have more prestigious jobs. Moreover, we would like to study whether professionals have more prestigious jobs than blue and white collar workers.

\newpage
\begin{enumerate}
	
	\item [(a)]
	Create a new variable \texttt{professional} by recoding the variable \texttt{type} so that professionals are coded as $1$, and blue and white collar workers are coded as $0$ (Hint: \texttt{ifelse}).
	
		\lstinputlisting[language=R, firstline=40, lastline = 40]{PS04.R}\\\\
		
	\vspace{1cm}
	
	
	\item [(b)]
	Run a linear model with \texttt{prestige} as an outcome and \texttt{income}, \texttt{professional}, and the interaction of the two as predictors (Note: this is a continuous $\times$ dummy interaction.)
			\lstinputlisting[language=R, firstline=44, lastline = 44]{PS04.R}\\\\
	\begin{table}[!htbp] \centering 
		\caption{Level of prestige based on income, professionals, and their interaction} 
		\label{} 
		\begin{tabular}{@{\extracolsep{5pt}}lc} 
			\\[-1.8ex]\hline 
			\hline \\[-1.8ex] 
			& \multicolumn{1}{c}{\textit{Dependent variable:}} \\ 
			\cline{2-2} 
			\\[-1.8ex] & prestige \\ 
			& Coefficients \\ 
			\hline \\[-1.8ex] 
			income & 0.003$^{***}$ \\ 
			& (0.0005) \\ 
			& \\ 
			professional & 37.781$^{***}$ \\ 
			& (4.248) \\ 
			& \\ 
			income:professional & $-$0.002$^{***}$ \\ 
			& (0.001) \\ 
			& \\ 
			Constant & 21.142$^{***}$ \\ 
			& (2.804) \\ 
			& \\ 
			\hline \\[-1.8ex] 
			Observations & 98 \\ 
			R$^{2}$ & 0.787 \\ 
			Adjusted R$^{2}$ & 0.780 \\ 
			Residual Std. Error & 8.012 (df = 94) \\ 
			F Statistic & 115.878$^{***}$ (df = 3; 94) \\ 
			\hline 
			\hline \\[-1.8ex] 
			\textit{Note:}  & \multicolumn{1}{r}{$^{*}$p$<$0.1; $^{**}$p$<$0.05; $^{***}$p$<$0.01} \\ 
		\end{tabular} 
	\end{table} 
	
	\vspace{6cm}
	\item [(c)]
	Write the prediction equation based on the result.
	\begin{equation}
		\widehat{\text{prestige}}  = 21.14 + 0.003\times \text{income} + 37.78\times\text{professional} - 0.002 \times\text{income * professional} 			
	\end{equation}
	\vspace{0.5 cm}
	
	\item [(d)]
	Interpret the coefficient for \texttt{income}.
	\begin{description}
		\item When professional categorgy equals zero (blue collar or white collar professions), a one unit increase in income is associated with a 0.003 increase in prestige. 
	\end{description}
	\vspace{1 cm}	
	\item [(e)]
	Interpret the coefficient for \texttt{professional}.
	\begin{description}
		\item When income equals zero, a one unit increase in professional is associated with a 37.78 increase in prestige. 
	\end{description}
	
	\vspace{0.5 cm}} \
	\item [(f)]
	What is the effect of a \$1,000 increase in income on prestige score for professional occupations? In other words, we are interested in the marginal effect of income when the variable \texttt{professional} takes the value of $1$. Calculate the change in $\hat{y}$ associated with a \$1,000 increase in income based on your answer for (c).
	\begin{description}
		\item When calculating the marginal effect, we are looking directly at the variables that have to do with income because we are looking at the effect of income on prestige score given the parameters of professional occupation and 1,000 dollar increase. Therefore, we take the coefficents that both include income, add them together, and multiply by the variable of professional when it equals 1 to get the marginal effect of income. Then to get the $\hat{y}$, we need to multiply this margin effect on the income parameter to see how much the 1,000 dollar increase in income of professional occupations impacts the prestige score.
	\end{description}
\begin{gather*}
	\text{marginal effect}  =  B1 + B3 * \text{professional} \\
	\text{marginal effect}  =  0.003 + (-0.002) * 1 \\
	\text{marginal effect}  =  0.001 \\
	\hat{y} = \text{marginal effect} * 1,000 \\
	\hat{y} = 0.001 * 1,000 = 1
\end{gather*}
\begin{description}
	\item An 1,000 dollar increase in income, increases a professional's prestige score by 1 point. 
\end{description}
	\vspace{2cm}
	
	\item [(g)]
	What is the effect of changing one's occupations from non-professional to professional when her income is \$6,000? We are interested in the marginal effect of professional jobs when the variable \texttt{income} takes the value of $6,000$. Calculate the change in $\hat{y}$ based on your answer for (c).
	\begin{description}
		\item For this scenerio, we are looking at the change from non-professional to professional on one's prestige score when her income is 6,000 dollars. So we take the coefficents that include the variable profession and multiple it by the income of 6,000 dollars to get the difference in prestige score from a non-professional occupation to a professional one.  
	\end{description}
	\begin{gather*}
		\text{marginal effect} = B2 + B3 *\text{income}\\
		\text{marginal effect} = 37.78 + (-0.002) * 6,000 = 25.78
	\end{gather*}
	
	\begin{description}
		\item Changing occupations from non-professional to professional when her income is 6,000 dollars, increases prestige score by 25.78 points. 
	\end{description}
\end{enumerate}

\newpage

\section*{Question 2: Political Science}
\vspace{.25cm}
\noindent 	Researchers are interested in learning the effect of all of those yard signs on voting preferences.\footnote{Donald P. Green, Jonathan	S. Krasno, Alexander Coppock, Benjamin D. Farrer,	Brandon Lenoir, Joshua N. Zingher. 2016. ``The effects of lawn signs on vote outcomes: Results from four randomized field experiments.'' Electoral Studies 41: 143-150. } Working with a campaign in Fairfax County, Virginia, 131 precincts were randomly divided into a treatment and control group. In 30 precincts, signs were posted around the precinct that read, ``For Sale: Terry McAuliffe. Don't Sellout Virgina on November 5.'' \\

Below is the result of a regression with two variables and a constant.  The dependent variable is the proportion of the vote that went to McAuliff's opponent Ken Cuccinelli. The first variable indicates whether a precinct was randomly assigned to have the sign against McAuliffe posted. The second variable indicates
a precinct that was adjacent to a precinct in the treatment group (since people in those precincts might be exposed to the signs).  \\

\vspace{.5cm}
\begin{table}[!htbp]
	\centering 
	\textbf{Impact of lawn signs on vote share}\\
	\begin{tabular}{@{\extracolsep{5pt}}lccc} 
		\\[-1.8ex] 
		\hline \\[-1.8ex]
		Precinct assigned lawn signs  (n=30)  & 0.042\\
		& (0.016) \\
		Precinct adjacent to lawn signs (n=76) & 0.042 \\
		&  (0.013) \\
		Constant  & 0.302\\
		& (0.011)
		\\
		\hline \\
	\end{tabular}\\
	\footnotesize{\textit{Notes:} $R^2$=0.094, N=131}
\end{table}

\vspace{.5cm}
\begin{enumerate}
	\item [(a)] Use the results from a linear regression to determine whether having these yard signs in a precinct affects vote share (e.g., conduct a hypothesis test with $\alpha = .05$).
	\vspace{0.5 cm}
\begin{align*}
	H_0 &: \beta_\text{assigned} = 0 \\
	H_1 &: \beta_\text{assigned} \ne 0
\end{align*}
\begin{description}
	\item The null hypothesis states that these yard signs in a precinct affects the vote share of Ken Cuccinelli. The alternative states that these signs do affect the vote share. 
\end{description}
		\lstinputlisting[language=R, firstline=63, lastline = 74]{PS04.R}\\
\begin{description}
	\item The p-value of 0.0097 is smaller than the alpha of 0.05 ,therefore, we reject the null hypothesis that yard signs in a precinct have no affect on the vote share of Ken Cuccinelli. 
\end{description}		
	\item [(b)]  Use the results to determine whether being
	next to precincts with these yard signs affects vote
	share (e.g., conduct a hypothesis test with $\alpha = .05$).
	\begin{align*}
		H_0 &: \beta_\text{adjacent}= 0 \\
		H_1 &: \beta_\text{adjacent}\ne 0
	\end{align*}
	\begin{description}
		\item The null hypothesis states that being next to precincts with these yard signs has no affects on Ken Cuccinelli's vote share. 
	\end{description}
	\lstinputlisting[language=R, firstline=77, lastline = 80]{PS04.R}\\
	\begin{description}
		\item The p-value of  0.0016 is less than alpha 0.05, therefore, we reject the null hypothesis that being next to precincts with these yard signs has no affects on Ken Cuccinelli's vote share. 
	\end{description}
	\vspace{.5cm}
	\item [(c)] Interpret the coefficient for the constant term substantively.
	\begin{description}
		\item On average, 0.302 is the predicted value of proportion of the vote that went to Ken Cuccinelli when there are no signs against McAuliffe in a precinct, therefore, also no exposure to adjacent precincts. 
	\end{description}
		\vspace{.5cm}
	\item [(d)] Evaluate the model fit for this regression.  What does this	tell us about the importance of yard signs versus other factors that are not modeled?
\begin{align*}
	H_0 &: \beta_\text{assigned} = \beta_\text{adjacent} = 0 \\
	H_1 &: \text{At least one } \beta_\text{assigned}  \text{or} \beta_\text{adjacent}\ne 0
\end{align*}
	\lstinputlisting[language=R, firstline=83, lastline = 86]{PS04.R}\\
\begin{description}
	\item The p-value of 0.0018 is less than alpha 0.05. Therefore, we reject the null hypothesis that these signs in assigned and adjacent precincts have no affect on Ken Cuccinelli's vote share. 
\end{description}
\end{enumerate}  


\end{document}
